\subsubsection{The Horizon Problem}
As we saw, the comonving distance to the last scattering surface is around \(\qty{46.6}{\giga\lightyear}\), but we can also look at the comoving distance for the maximum distance light could have travelled from the Big Bang to the time of recombination. This is given by
\begin{equation}
    x_{\max} = \qty{0.06}{} ct_0
\end{equation}
Which is very small compared to the distance to the last scattering surface. This raises the question of how different regions of the CMB which are very far apart compared to \(x_{\max}\), meaning they did not have time to communicate, have the same temperature to such a high degree of accuracy. This is known as the horizon problem, and is one of the main motivations for the inflationary model. Another analysis that can be done on the CMB is to look at the angular size of the fluctuations. This can be done by decomposing the temperature fluctuations into spherical harmonics and the resulting graph of the mean square of the fluctuations as a fucntion of the multipole moment \(l\) shows peaks at certain values of \(l\). These peaks point at a characteristic length scale, which corresponds to the sound horizon at recombination. This is the maximum distance that a sound wave could have travelled from the Big Bang to the time of recombination, and is given by
\begin{equation}
    \frac{1}{2}\lambda_{\max} = c_s t_{rec}
\end{equation}
and if we insert numbers we find that \(\frac{1}{2}\lambda_{\max} \approx \qty{130}{\kilo\parsec}\). This gives us a standard ruler to measure distances in the universe. The angle that this length subtends on the sky can be found using
\begin{equation}
    \dl p = \frac{l}{\delta \theta} (1+z)
\end{equation}
if we insert all the numbers from theory we find a value of \(\delta \theta \approx \qty{0.7}{\degree}\) which matches the observations very well. This is another strong piece of evidence for the Big Bang model. The development of this equation as we recall assumed a flat universe, and since the measurement matches the flat universe prediction very well, this is another strong evidence that the universe is flat. If we use the same equation with \(\Omega_0\) we find that its value is
\begin{equation}
    \Omega_0 = 1.02 \pm 0.01
\end{equation}
\section{Spiral Galaxies In 10 Minutes}
Spiral galaxies are composed of three main components: a disk of stars and gas, a bulge in the center which is more spherical, and a halo which is a more diffuse spherical component. The halo contains elements called globular clusters which are very old star clusters. In the middle of the bulge there is a supermassive black hole, and we are located about \(\qty{8.5}{\kilo\parsec}\) from the center of the galaxy in one of the spiral arms. Another region beyond the halo is the dark matter halo which extends much further out. If we look at a figure of the density of stars or gas as a function of the radius from the center of the galaxy, starting at a certain radius the density starts to drop off exponentially. If we integrate over the plot we'll get a figure of the mass of the galaxy as a function of radius which will increase linearly until that radius then stay constant. If we then look at the rotation of the galaxy using circular motion
\begin{equation}
    \frac{GM(r)}{r^2} = \frac{v^2}{r}
\end{equation}
we find that the velocity should drop off as
\begin{equation}
    v  \propto r^{-1/2}
\end{equation}
However, if we look at the actual rotation curve of the galaxy we find that the velocity stays mostly constant as we go further out. This depends on the specific galaxy but generally the rotation curves are flat. We can use these velocities to find an experimental graph of the density and mass as a function of radius. This discrepancy between the expected and observed rotation curves is one of the main pieces of evidence for dark matter.